\subsection{\texorpdfstring{Computation of the modules \(\operatorname{Tor}^{\mathbf{A}}(\mathbf{P},\mathbf{M})\) and \(\operatorname{Ext}_{\mathbf{A}}(\mathbf{M},\mathbf{N})\)}{Computation of the modules \textbackslash operatorname\{Tor\}\^{}\{\textbackslash mathbf\{A\}\}(\textbackslash mathbf\{P\},\textbackslash mathbf\{M\}) and \textbackslash operatorname\{Ext\}\_\{\textbackslash mathbf\{A\}\}(\textbackslash mathbf\{M\},\textbackslash mathbf\{N\})}}

Let M, N be left A-modules, P a right A-module. Let moreover \(a: R \rightarrow M\) be a left resolution of M, \(b: S \rightarrow P\) a left resolution of P, and \(c: N \rightarrow E\) a right resolution of N.

By X, p.~49, prop. 3 and 3 bis, there exist morphisms of complexes \(\alpha: L(M) \to R\) , \(\beta: L(P) \to S\) , \(\gamma: E \to I(N)\) such that \(a \circ \alpha = p_{M}\) , \(b \circ \beta = p_{P}\) , \(\gamma \circ c = e_{N}\) ; and the homotopy classes of \(\alpha\) , \(\beta\) , \(\gamma\) depend only on the given resolutions. By X, p.~64, prop. 3 and X, p.~84, prop. 3, the homotopy classes of the morphisms

\[
\beta \otimes \alpha : L (P) \otimes_ {A} L (M) \rightarrow S \otimes_ {A} R,
\]

\[
\operatorname{Homgr} _ {\mathbf {A}} (\alpha , \gamma): \operatorname{Homgr} _ {\mathbf {A}} (\mathrm{R}, \mathrm{E}) \rightarrow \operatorname{Homgr} _ {\mathbf {A}} (\mathrm{L} (\mathrm{M}), \mathrm{I} (\mathrm{N}))
\]

depend only on the given resolutions, whence, by passing to homology, graded k-homomorphisms of degree 0

\[
\psi (\mathrm{S}, \mathrm{R}): \operatorname{Tor} ^ {\mathrm{A}} (\mathrm{P}, \mathrm{M}) \rightarrow \mathrm{H} (\mathrm{S} \otimes_ {\mathrm{A}} \mathrm{R}),
\]

\[
\varphi (\mathrm{R}, \mathrm{E}): \mathrm{H} \left(\operatorname{Homgr} _ {\mathrm{A}} (\mathrm{R}, \mathrm{E})\right)\rightarrow \operatorname{Ext} _ {\mathrm{A}} (\mathrm{M}, \mathrm{N}),
\]

independent of the choices of α, β, γ.

For example, taking for \(a\) , \(b\) , \(c\) the identity maps of M, P, N respectively, one obtains the homomorphisms \(\psi(\mathbf{P}, \mathbf{M}) : \operatorname{Tor}^{\mathbf{A}}(\mathbf{P}, \mathbf{M}) \to \mathbf{P} \otimes_{\mathbf{A}} \mathbf{M}\) and \(\varphi(\mathbf{M}, \mathbf{N}) : \operatorname{Hom}_{\mathbf{A}}(\mathbf{M}, \mathbf{N}) \to \operatorname{Ext}_{\mathbf{A}}(\mathbf{M}, \mathbf{N})\) introduced in X, p.~68, remark 2) and X, p.~87, remark 2).

THEOREM 1. --- a) If one of the resolutions R or S is flat, then \(\psi(S, R)\) is an isomorphism of \(k\) -graded modules.

\begin{enumerate}
\def\labelenumi{\alph{enumi})}
\setcounter{enumi}{1}
\item
  If R is projective or if E is injective, \(\varphi(\mathbf{R},\mathbf{E})\) is an isomorphism of \(k\) -graded modules.
\item
  Suppose for example R is flat, and choose \(\alpha\) and \(\beta\) as above. The homomorphism \(\beta \otimes \alpha\) is the composite of the morphisms
\end{enumerate}

\[
\mathrm{L} (\mathrm{P}) \otimes_ {\mathrm{A}} \mathrm{L} (\mathrm{M}) \xrightarrow {l _ {\mathrm{L} (\mathrm{P})} \otimes \alpha} \mathrm{L} (\mathrm{P}) \otimes_ {\mathrm{A}} \mathrm{R} \xrightarrow {\beta \otimes 1 _ {\mathrm{R}}} \mathrm{S} \otimes_ {\mathrm{A}} \mathrm{R}.
\]

Since L(P) (resp. R) is flat and \(\alpha\) (resp. \(\beta\) ) is a homology isomorphism, \(1_{L(P)} \otimes \alpha\) (resp. \(\beta \otimes 1_R\) ) is one, by prop. 4 of X, p.~67. Hence \(\beta \otimes \alpha\) is a homology isomorphism and \(\psi(S, R) = H(\beta \otimes \alpha)\) is bijective.

\begin{enumerate}
\def\labelenumi{\alph{enumi})}
\setcounter{enumi}{1}
\tightlist
\item
  One argues similarly, using prop. 4 of X, p.~86.
\end{enumerate}

COROLLARY. --- If R is a flat resolution of M, the homomorphism

\[
\psi (\mathrm{P}, \mathrm{R}): \operatorname{Tor} ^ {\mathbf {A}} (\mathrm{P}, \mathrm{M}) \rightarrow \mathrm{H} (\mathrm{P} \otimes_ {\mathbf {A}} \mathrm{R})
\]

is bijective. If R is a projective resolution of M, the homomorphism

\[
\varphi (\mathrm{R}, \mathrm{N}): \mathrm{H} \left(\operatorname{Homgr} _ {\mathrm{A}} (\mathrm{R}, \mathrm{N})\right)\rightarrow \operatorname{Ext} _ {\mathrm{A}} (\mathrm{M}, \mathrm{N})
\]

is bijective. If E is an injective resolution of N, the homomorphism

\[
\varphi (\mathbf {M}, \mathbf {E}): \mathrm{H} \left(\operatorname{Homgr} _ {\mathbf {A}} (\mathbf {M}, \mathbf {E})\right)\rightarrow \operatorname{Ext} _ {\mathbf {A}} (\mathbf {M}, \mathbf {N})
\]

is bijective.

Remark. --- The diagram of k-modules

\[
\begin{array}{c} \operatorname{Tor} ^ {\mathbf {A}} (\mathrm{P}, \mathrm{M}) \xrightarrow {\psi (\mathrm{S} , \mathrm{R})} \mathrm{H} (\mathrm{S} \otimes_ {\mathbf {A}} \mathrm{R}) \\ \sigma_ {\mathrm{P}, \mathrm{M}} \Biggl \downarrow \\ \operatorname{Tor} ^ {\mathbf {A} ^ {\circ}} (\mathrm{M} ^ {\circ}, \mathrm{P} ^ {\circ}) \xrightarrow {\psi (\mathrm{R} ^ {\circ} , \mathrm{S} ^ {\circ})} \mathrm{H} (\mathrm{R} ^ {\circ} \otimes_ {\mathbf {A} ^ {\circ}} \mathrm{S} ^ {\circ}), \end{array}
\]

where \(\sigma_{\mathrm{P,M}}\) and \(\sigma(\mathrm{S},\mathrm{R})\) are the commutation isomorphisms (X, p.~71 and 63), is commutative: it is indeed deduced, by passing to homology, from the commutative diagram of complexes

\[
\begin{array}{c c} \mathrm{L} (\mathrm{P}) \otimes_ {\mathrm{A}} \mathrm{L} (\mathrm{M}) & \xrightarrow {\beta \otimes \alpha} \mathrm{S} \otimes_ {\mathrm{A}} \mathrm{R} \\ \sigma (\mathrm{L} (\mathrm{P}), \mathrm{L} (\mathrm{M})) \Big \downarrow & \Big \downarrow \sigma (\mathrm{S}, \mathrm{R}) \\ \mathrm{L} (\mathrm{M} ^ {\circ}) \otimes_ {\mathrm{A} ^ {\circ}} \mathrm{L} (\mathrm{P} ^ {\circ}) & \xrightarrow {\alpha \otimes \beta} \mathrm{R} ^ {\circ} \otimes_ {\mathrm{A} ^ {\circ}} \mathrm{S} ^ {\circ} \end{array}
\]

(“the morphisms \(\psi\) are compatible with the commutation isomorphisms”).

Similarly, let \(a_{1}: \mathbb{R}_{1} \to \mathbb{M}\) , \(b_{1}: \mathbb{S}_{1} \to \mathbb{P}\) , \(c_{1}: \mathbb{N} \to \mathbb{E}_{1}\) be morphisms of complexes, where \(\mathbb{R}_{1}\) and \(\mathbb{S}_{1}\) are projective and zero on the right, and \(\mathbb{E}_{1}\) injective and zero on the left. By X, p.~49, prop. 3 and 3 bis, there exist morphisms of complexes

\[
\alpha_ {1}: \mathrm{R} _ {1} \rightarrow \mathrm{L} (\mathrm{M}), \quad \beta_ {1}: \mathrm{S} _ {1} \rightarrow \mathrm{L} (\mathrm{P}), \quad \gamma_ {1}: \mathrm{I} (\mathrm{N}) \rightarrow \mathrm{E} _ {1}
\]

such that \(p_{\mathrm{M}} \circ \alpha_{1} = a_{1}, p_{\mathrm{P}} \circ \beta_{1} = b_{1}, \gamma_{1} \circ e_{\mathrm{N}} = c_{1}\) , whence morphisms of complexes

\[
\beta_ {1} \otimes \alpha_ {1}: S _ {1} \otimes_ {A} R _ {1} \rightarrow L (P) \otimes_ {A} L (M),
\]

\[
\operatorname{Homgr} _ {\mathbf {A}} \left(\alpha_ {1}, \gamma_ {1}\right): \operatorname{Homgr} _ {\mathbf {A}} \left(\mathrm{L} (\mathrm{M}), \mathrm{I} (\mathrm{N})\right)\rightarrow \operatorname{Homgr} _ {\mathbf {A}} \left(\mathrm{R} _ {1}, \mathrm{E} _ {1}\right),
\]

and, by passing to homology, graded k-linear maps of degree 0:

\[
\begin{array}{l} \psi^ {\prime} (\mathbf {S} _ {1}, \mathbf {R} _ {1}): \mathrm{H} (\mathbf {S} _ {1} \otimes_ {\mathbf {A}} \mathbf {R} _ {1}) \to \operatorname{Tor} ^ {\mathbf {A}} (\mathbf {P}, \mathbf {M}) \\ \varphi^ {\prime} (\mathbf {R} _ {1}, \mathbf {E} _ {1}): \operatorname{Ext} _ {\mathbf {A}} (\mathbf {M}, \mathbf {N}) \to \mathrm{H} (\operatorname{Homgr} _ {\mathbf {A}} (\mathbf {R} _ {1}, \mathbf {E} _ {1})) \end{array}
\]

which one verifies as above are independent of the choice of \(\alpha_{1}\) , \(\beta_{1}\) , \(\gamma_{1}\) .

PROPOSITION 1. --- If \(a_{1}\) , \(b_{1}\) , \(c_{1}\) are homology isomorphisms, \(\psi^{\prime}(\mathbf{S}_{1},\mathbf{R}_{1})\) and \(\varphi^{\prime}(\mathbf{R}_{1},\mathbf{E}_{1})\) are the inverse bijections of the bijections \(\psi(\mathbf{S}_{1},\mathbf{R}_{1})\) and \(\varphi(\mathbf{R}_{1},\mathbf{E}_{1})\) respectively. Indeed, \(f = (\beta \otimes \alpha) \circ (\beta_{1} \otimes \alpha_{1})\) is a morphism of the complex \(\mathbf{S}_{1} \otimes_{\mathbf{A}} \mathbf{R}_{1}\) into itself, and one has \((b_{1} \circ \alpha_{1}) \circ f = f\) . By X, p.~49, prop. 3 and 3 bis, \(f\) is a homotopy equivalence, hence H(f) = 1 and

\[
\psi \left(\mathrm{S} _ {1}, \mathrm{R} _ {1}\right) \circ \psi^ {\prime} \left(\mathrm{S} _ {1}, \mathrm{R} _ {1}\right) = \mathrm{H} (\beta \otimes \alpha) \circ \mathrm{H} \left(\beta_ {1} \otimes \alpha_ {1}\right) = \mathrm{H} (f) = 1;
\]

likewise \(\psi'(S_1, R_1) \circ \psi(S_1, R_1) = 1\) . One argues analogously for the maps \(\varphi\) and \(\varphi'\) .

Examples. --- 1) Let \(a\) be an element of A such that the map \(\varphi : x \mapsto xa\) from A into itself be injective ( \(a\) is not a right zero divisor). Using the resolution

\[
0 \rightarrow \mathrm{A} _ {s} \xrightarrow {\varphi} \mathrm{A} _ {s} \rightarrow \mathrm{A} / \mathrm{A} a \rightarrow 0
\]

we see that for every right A-module M, we have

\[
\operatorname{Tor} _ {i} ^ {\mathbf {A}} (\mathbf {M}, \mathbf {A} / \mathbf {A} a) = 0 \quad \text { for } \quad i > 1
\]

and that the \(k\) -module \(\operatorname{Tor}_1^{\mathbf{A}}(\mathbf{M}, \mathbf{A} / \mathbf{A}a)\) is isomorphic to \(\operatorname{Ker}(a_{\mathbf{M}})\) .

Similarly, for every left A-module M, we have

\[
\operatorname{Ext} _ {\mathbf {A}} ^ {i} (\mathbf {A} / \mathbf {A} a, \mathbf {M}) = 0 \quad \text { for } \quad i > 1
\]

and the \(k\) -module \(\operatorname{Ext}_{\mathbf{A}}^{1}(\mathbf{A}/\mathbf{A}a, \mathbf{M})\) is isomorphic to \(\mathbf{M}/a\mathbf{M}\) .

\begin{enumerate}
\def\labelenumi{\arabic{enumi})}
\setcounter{enumi}{1}
\tightlist
\item
  Suppose A is an integral domain; let K be the field of fractions of A and M an A-module. Using the flat resolution
\end{enumerate}

\[
0 \rightarrow \mathrm{A} \rightarrow \mathrm{K} \rightarrow \mathrm{K} / \mathrm{A} \rightarrow 0
\]

(X, p.~9, example 5), we see that \(\operatorname{Tor}_{i}^{\mathbf{A}}(\mathbf{K}/\mathbf{A}, \mathbf{M}) = 0\) for \(i > 1\) ; moreover, taking into account II, p.~116, prop. 26, (ii), the A-module \(\operatorname{Tor}_{1}^{\mathbf{A}}(\mathbf{K}/\mathbf{A}, \mathbf{M})\) is isomorphic to the torsion submodule of M.

\begin{enumerate}
\def\labelenumi{\arabic{enumi})}
\setcounter{enumi}{2}
\tightlist
\item
  Suppose that A is a Noetherian local ring, denote its maximal ideal by m, and set \(\kappa = \mathbf{A} / \mathfrak{m}\) Let M be a finitely generated A-module and P a minimal projective resolution of M (X, p.~54). For every \(n \geqslant 0\) , the \(\kappa\) vector spaces \(\operatorname{Tor}_n^{\mathbf{A}}(\kappa, \mathbf{M})\) and \(\operatorname{Ext}_{\mathbf{A}}^{n}(\mathbf{M}, \kappa)\) are finite-dimensional, of dimension equal to the rank of the free A-module \(\mathbf{P}_n\) ; indeed, the complexes \(\kappa \otimes_{\mathbf{A}} \mathbf{P}\) and \(\operatorname{Homgr}_{\mathbf{A}}(\mathbf{P}, \kappa)\) have zero differential.
\end{enumerate}

